\section{A signed-tensor overlap with a minimized tag budget}
\label{ten:compressed}
The argument here minimizes an entropy cost over two signed-spin partitions and a freely labeled remainder. The resulting budget can be negative, so the finite-size and sparse parts of the induction include an explicit additive allowance. The overlap estimate retains its rank dependence.
\smallskip
The notation below is local; the sweep and trace conventions remain those of Section~\ref{sec:prelim}.

\subsection{The budget and the moment statement}
For a list \(v\) of part sizes write
\(h(v)=\sum v_i\log(\sum_j v_j/v_i)\), omitting zeros. Whenever
\(M>0\), we use the continuous convention \(x\log(M/x)=0\) at \(x=0\).
Define a budget
\[
L_n(\lambda)=\min\{ b h(\mu,\xi)+ c t\log n-t\log(n/t)\}.
\]
Here the diagrams \(\mu,\xi\) have total size \(n-t\), and \(\lambda\) occurs in the representation induced by \(\mu,\xi^\top\) on disjoint blocks with the remaining \(t\) sites freely labeled (a regular factor). Use absolute positive constants \(b\) sufficiently small, \(c\) sufficiently small compared with \(b\).

\begin{proposition}[Moment estimate with a minimized tag budget]
\label{ten:compressed:main}
For sufficiently small fixed \(b>0\) and then sufficiently small fixed
\(c>0\), there are a constant \(C\) and exponents \(p_d\ge2\) with
\(\sup_{d\ge1}p_d<\infty\) such that, for every \(d\ge1\) and
every \(\lambda\vdash n=2^d\),
\begin{equation}
 D_\lambda {\rm tr}|\rho_\lambda(B_n)|^{p_d}\le
 \exp\{L_n(\lambda)+C n^{1-c/2}\}.
 \label{ten:compressed:eq:1}
\end{equation}
\end{proposition}

\subsection{Three budget properties}
Here are useful budget facts.
\begin{itemize}
\item \(L_n\le\frac12\log D_\lambda\).
Peel off successively the longer of the first row and first column.
Write $\mathbf r=(r_i)$ for the row lengths peeled, $\mathbf c=(c_j)$ for the column
lengths peeled, and $v$ for all these lengths in their peeling order.
A row peel leaves an interlacing partition, so its reversal adds
a horizontal strip; a column peel reverses to a vertical strip.
Pieri's rule therefore gives a trivial factor for a row peel and
an alternating factor for a column peel. Put
$n_+=\sum_i r_i$, $n_-=\sum_j c_j$, and
\[
 M_{\mathbf r}=\operatorname{Ind}_{\prod_i S_{r_i}}^{S_{n_+}}\mathbf1,
 \qquad
 M_{\mathbf c}=\operatorname{Ind}_{\prod_j S_{c_j}}^{S_{n_-}}\mathbf1.
\]
Induction by stages, grouping the row and column factors, shows that
$V_\lambda$ occurs in
\[
 \operatorname{Ind}_{S_{n_+}\times S_{n_-}}^{S_n}
   \bigl(M_{\mathbf r}\otimes
         (\operatorname{sgn}_{n_-}\otimes M_{\mathbf c})\bigr).
\]
Decompose $M_{\mathbf r}$ and $M_{\mathbf c}$ into irreducibles and choose constituents
$V_\mu$ and $V_\xi$ whose corresponding induced summand still
contains $V_\lambda$. Since
$\operatorname{sgn}_{n_-}\otimes V_\xi\simeq V_{\xi^\top}$,
this is an admissible pair for the budget with $t=0$.
An empty list uses the trivial representation of $S_0$.
Young's rule~\cite[14.1, pp.~51--52]{James1978}
says that $\mu$ and $\xi$ dominate the decreasing
rearrangements of $\mathbf r$ and $\mathbf c$, respectively. Convexity of
$x\log x$, applied separately at the fixed totals $n_+,n_-$, gives
\[
 \begin{aligned}
 h(\mu,\xi)
 &=n\log n-\sum_i\mu_i\log\mu_i-\sum_j\xi_j\log\xi_j\\
 &\le n\log n-\sum_i r_i\log r_i-\sum_j c_j\log c_j
 =h(v).
 \end{aligned}
\]
Thus $L_n\le b h(v)$.

To bound $h(v)$, write \(\ell=v_1,\ s=n-\ell\).
If $s=0$, then $h(v)=0$.
The hook product on the first removed line is at most
\(\ell!\exp(s)\), by adding the \(s\) excess lengths; on every
later one it is at most \((2v_i)^{v_i}\). Stirling or factorial
estimates thus give \(\log D_\lambda\ge h(v)-O(s)\) (the ratio
\(n!/\ell!\) has log at least
\(n\log n-\ell\log\ell-O(s)\)). Also
\(\log D_\lambda=\Omega(s)\) for growing \(n\): if \(\ell\)
is a sufficiently small fraction of \(n\) this follows directly;
otherwise keep the first line and a subdiagram below (or beside)
it of size \(u=\min(s,\lfloor\ell/3\rfloor)\); interleavings of
its fixed tableau with the line respecting the two-list ballot
condition already give exponentially many in \(s\). Bounded
small sizes can equivalently be handled separately. Hence
$h(v)\le C_0\log D_\lambda$ uniformly for an absolute $C_0$, and fixing
$b\le1/(2C_0)$ proves the claimed budget bound.
\item Every minimizing choice has at most \(q=\lceil n^{1/16}\rceil\) parts in each of \(\mu,\xi\), for large \(n\). Otherwise removing one end box from a smallest part and making it a tag (increasing \(t\)) saves \(b((1/16)\log n-O(1))\) while costing at most \(c\log n+1\). Branching permits this choice.
\item \(L_n\ge -O(n^{1-c/2})\), and at the minimum \(t=O(n^{1-c/2}+\log D_\lambda/\log n)\).
\end{itemize}

For the last property, isolate the raw tag term on the real interval \(0\le t\le n\):
\[
 \begin{aligned}
 \phi_n(0)&=0,\\
 \phi_n(t)&=ct\log n-t\log(n/t)\\
          &=t\log(t/n^{1-c})\quad(0<t\le n).
 \end{aligned}
\]
Its derivative is \(\log(t/n^{1-c})+1\), so its minimum occurs at
\(t=n^{1-c}/e\) and equals \(-n^{1-c}/e\). For
\(t\ge n^{1-c/2}\), we instead have
\(\phi_n(t)\ge(c/2)t\log n\). Since \(bh\ge0\), the first bound gives
\(L_n\ge-n^{1-c}/e\ge-n^{1-c/2}/e\). At a minimizing choice above the
second threshold,
\((c/2)t\log n\le L_n\le\frac12\log D_\lambda\), so
\(t\le\log D_\lambda/(c\log n)\); below that threshold,
\(t<n^{1-c/2}\). These are the two assertions in the last bullet.



In the dense-level case \(n-\lambda_1\ge n^{1-\eta}\), where \(\eta>0\) is fixed sufficiently small compared with \(c\), only \(\ell(\lambda)\le n/2\) matters: taller diagrams have no invariant vector for the first matching (Young's rule). Here \(\log D_\lambda\gtrsim n^{1-\eta}\) by the preceding estimates. Consequently multiplicative losses \(\exp O(t+n^{1-c/6})\) in proving \eqref{ten:compressed:eq:1} at such levels can be absorbed (even to give the bound without the \(C n^{1-c/2}\)) by raising the exponent by a factor \(1+O(1/d)\), since the bound itself then forces each singular value to power \(p_d\) to be at most \(\exp(-\Omega(\log D_\lambda))\).

\subsection{Sparse mixed-particle paths}
First we record the sparse-level estimate for \(1\le k=n-\lambda_1<n^{1-\eta}\):
\begin{equation}
 \|\rho_\lambda(B_n)\|\le n^{-a k}
 \label{ten:compressed:eq:2}
\end{equation}
for large \(d\), some \(a(\eta)>0\). Work on the module of injections of \(k\) tracked labels. Lower-coordinate kernels (depending on fewer than all the labels of the output) kill the isotype in question, by branching. Analyze \(B_n B_n^* B_n\); in the middle sweep replace the joint kernel by the alternating inclusion-exclusion sum over subsets \(I\) of labels, in each term updating \(I\) together and each of its complement by private independent randomness. Use only distinct endpoints. Proper subsets give lower-coordinate kernels since a single particle finishes exactly uniformly. Given start and finish the paths are forced within the sweep. Draw their contact graph (using a common switch location at the same time). If some label is isolated the inclusion-exclusion entries cancel exactly (factorization on disjoint switches), so restrict every summand to no isolates. Each resulting nonnegative matrix, including the flanking sweeps, has column sums at most 1 (without restriction the independent-group updates preserve counting measure on their product space).

Its row sums are at most \(n^{-\Omega_\eta(k)}\), as follows. After the first sweep the tracked occupied sites have joint occupation upper bounds \((k/n)^r\) for every set of \(r\) sites. Indeed joint occupancy bounded by products of one-site probabilities persists from deterministic sets under fair transpositions (average the products, using arithmetic-geometric mean when both sites occur).

Fix the original injection, the common subset \(I\), a time \(\theta\) in the mixed middle sweep, and an ordered list of targets
\(\mathbf y=(y_1,\ldots,y_r)\), where \(r\ge1\) and repetitions are allowed.
Let \(Z_i(\theta)\) be the position of label \(i\) after the first true sweep and the indicated middle layers. For
\(J\subseteq\{1,\ldots,r\}\), define the following sums over tracked labels:
\[
 \begin{aligned}
 N_J(\mathbf y)&=
 \sum_{\substack{i_1,\ldots,i_r\text{ distinct}\\
                  i_h\in I\ \Longleftrightarrow\ h\in J}}
       \prod_{h=1}^r \mathbf1\{Z_{i_h}(\theta)=y_h\},\\
 N(\mathbf y)&=\sum_JN_J(\mathbf y).
 \end{aligned}
\]
Thus \(N\) counts ordered distinct labels, not necessarily distinct target sites.

Temporarily condition on \(\mathcal F_0\), the sigma-field generated by the random starting configuration of the middle sweep. The coordinates updated by time \(\theta\) partition the sites into blocks \(C\) of a common size \(v\). In block \(C\), let \(A_C,F_C\) be the initial common and private particle counts, and let \(j_C,\ell_C\) count target indices in \(J,J^c\). If two common targets coincide, then \(N_J=0\), because common particles remain at distinct sites. Otherwise the fair-switch occupation bound applies to the \(j_C\) common targets; the distinct private labels have independent uniform positions and contribute the falling factorial \((F_C)_{\ell_C}\). Conditional independence between blocks and between the common and private randomness gives
\[
 \begin{split}
 \mathbb E[N_J(\mathbf y)\mid\mathcal F_0]
 &\le \prod_C
      \frac{\mathbf1_{\{A_C\ge j_C\}}A_C^{j_C}(F_C)_{\ell_C}}
           {v^{j_C+\ell_C}}\\
 &\le \prod_C
      \frac{e^{j_C}(A_C+F_C)_{j_C+\ell_C}}
           {v^{j_C+\ell_C}}.
 \end{split}
\]
Private targets may repeat, and a common and private target may coincide. We use \((a)_0=a^0=1\), with infeasible falling factorials zero. For \(A\ge j\), the second inequality follows from
\(A^j/(A)_j\le j^j/j!\le e^j\) and
\((A)_j(F)_\ell\le(A+F)_{j+\ell}\), with the evident convention at \(j=0\).

Put \(r_C=j_C+\ell_C\). Averaging over the first true sweep now gives
\[
 \mathbb E\prod_C(A_C+F_C)_{r_C}
 \le \left(\prod_C(v)_{r_C}\right)(k/n)^r
 \le v^r(k/n)^r.
\]
Indeed the factorial product counts ordered distinct occupied starting sites within disjoint blocks, and each set of \(r\) such sites has joint occupation probability at most \((k/n)^r\). Averaging the conditional bound and summing over \(J\) therefore proves
\[
 \mathbb E N(\mathbf y)
 \le (1+e)^r(k/n)^r
 \le (2e k/n)^r.
\]
This final expectation includes the first true sweep and the middle randomness; it conditions only on the original injection and \(I\), with \(\theta\) and \(\mathbf y\) prescribed.

Divide the middle sweep into a constant (depending on \(\eta\)) number of pieces each varying a block of bits of volume at most \(n^{\eta/2}\). No isolates forces in some piece \(\Omega_\eta(k)\) particles sharing their fixed-bit classes with peers at its start, hence at least \(r=\Omega_\eta(k)\) disjoint pairs there (round with \(r\ge1\), excluding the impossible case). Counting ordered target pairs, applying the preceding occupation bound to their \(2r\) ordered endpoints, and dividing by \(r!\), probability at most
\[
 (n^{1+\eta/2})^r(2ek/n)^{2r}/r!.
\]
This proves the assertion. The matrix row-column sum bound and \(2^k\) terms now bound the cubed singular norm, proving \eqref{ten:compressed:eq:2}. Since \(D_\lambda\le n^k\), this suffices both for \eqref{ten:compressed:eq:1} and for stronger decay with a fixed exponent.

\subsection{The dense-level trace induction}
We detail the dense-level induction. Use the two halves of the bits to arrange the sites in a rectangle with side lengths \(m=2^{\lfloor d/2\rfloor},m'=n/m\). (Call the row length \(m\).) The sweep factors are products of independent smaller sweeps on rows and on columns. With \(p=\max(p_{\lfloor d/2\rfloor},p_{\lceil d/2\rceil})\), bound the left trace without \(D_\lambda\) by \({\rm tr}_\lambda AB\), where \(A,B\) are positive subgroup algebras on rows and columns respectively, absolute powers \(p\) as in the trace inequality. Each is a tensor product, with regular trace budgets \eqref{ten:compressed:eq:1} per isotype per line.

Use minimizing \(\mu,\xi,t\), choosing \(\nu\vdash t\) such that induction with the tag block of type \(\nu\) contains \(\lambda\); in particular
\[
 D_\lambda\le \binom{n}{t}e^{h(\mu,\xi)}D_\nu .
\]
Majorize the trace by using induction with \textbf{regular} tag block divided by \(D_\nu\), allowed by positivity of irrep traces. In expanding \(AB\), a trace contribution fixes all tagged sites individually. Since row and column lines intersect once, each factor permutation must fix them individually. Thus the upper bound is \(t!/D_\nu\) times a sum over tag sets \(T\) of size \(t\), of \({\rm tr}_{\mathcal H} A_0 B_0\). Here \(A_0,B_0\) are coefficient restrictions to row and column groups avoiding \(T\), and \(\mathcal H\) on the remaining grid sites is the representation induced by \(\mu,\xi^\top\). Restrictions are positive (compress regular matrices).

Write $r_i$ for row tag counts and $s_j$ for column counts. Let
$H$ be the product of the full row groups and $H_0$ its subgroup
fixing the tags pointwise; write $J,J_0$ for columns. A row type is
an irreducible representation $V_\gamma=\bigotimes_iV_{\gamma_i}$
of $H_0$, of dimension $D_\gamma=\prod_iD_{\gamma_i}$.
The element $A_0\in\mathbb C[H_0]$ acts by the same matrix on every
copy of this type. In the next display,
$\operatorname{tr}(A_0)_\gamma$ denotes the unnormalized trace on
one such copy. For each row type occurring on $\mathcal H$,
\begin{equation}
 D_\gamma {\rm tr}(A_0)_\gamma
 \le \frac{|H_0|}{|H|}
 \exp\{b\log D_\gamma+c t\log m-\sum_i r_i\log(m/r_i)
            +O(n^{1-c/6})\}.
 \label{ten:compressed:eq:3}
\end{equation}
The column estimate is identical. To prove the row estimate, apply
Lemma~\ref{lem:pointwise-positive-restriction} in each row, with
$G=S_m$, its subgroup $S_{m-r_i}$, and the positive local factor
$A_i$ of $A$. It gives
\[
 D_{\gamma_i}\operatorname{tr}_{\gamma_i}
       (E_{S_{m-r_i}}A_i)
 \le\frac1{(m)_{r_i}}
       \sum_{\alpha_i\supseteq\gamma_i}
       D_{\alpha_i}\operatorname{tr}_{\alpha_i}A_i.
\]
Thus only extending parent types contribute, and there is no
branching-multiplicity factor in this sum. Multiplying over rows
gives $\prod_i(m)_{r_i}^{-1}=|H_0|/|H|$.
The number of extension choices is at most
$\exp O(n^{3/4}\log n)$ by the partition bound.

Each $\gamma_i$ occurs in induction of a signed pair
$\delta,\epsilon^\top$, both of length at most $q$: restrict the
global signed tensor realization to the row and split its two parity
classes. Every extension $\alpha_i$ can then use this pair and
$r_i$ regular tags as a candidate in $L_m(\alpha_i)$.
The minimum is at most this candidate cost, so the child bound
\eqref{ten:compressed:eq:1} applies with
\begin{equation}
 h(\delta,\epsilon)\le\log D_{\gamma_i}+O(q^2\log n).
 \label{ten:compressed:eq:4}
\end{equation}
For clarity, any Littlewood-Richardson result \(\gamma_i\) of \(\delta,\epsilon^\top\) contains both diagrams. Outside their union it adds at most \(q^2\) boxes since that bounds their intersection. In the first \(q\) rows and columns where they overlap or interact (the \(q\times q\) corner) hook costs are polynomial per box; in the horizontal arm, for boxes of \(\delta\), downward hook contributions within first \(q\) rows cost at most \(O(q^2\log n)\) extra log over row factorials. Extra boxes outside \(\delta,\epsilon^\top\) also modify hooks by \(O(\log n)\) per box (harmonic sums along their lines); boxes of \(\epsilon^\top\) do not extend into the horizontal arm, but can lengthen rows there only via the corner, already within the row extent of \(\delta\) if that row has arm boxes. The vertical arm works symmetrically. Thus hook product is at most \(\prod\delta_i!\prod\epsilon_j!\exp O(q^2\log n)\), giving \eqref{ten:compressed:eq:4}.

The child allowances sum over the $n/m$ rows to
$C(n/m)m^{1-c/2}=O(n^{1-c/4})$.
Together with the type-entropy errors and extension counts, this
fits $O(n^{1-c/6})$ for the fixed small $c$. The remaining child
terms sum to $b\log D_\gamma+ct\log m-\sum_i r_i\log(m/r_i)$,
which proves \eqref{ten:compressed:eq:3} with constants independent
of the moment exponent.

\subsection{The signed-tensor overlap}
We next retain the ranks needed to combine the one-copy traces in
\eqref{ten:compressed:eq:3}. Fix a row type $V_\gamma$ of $H_0$
and a column type $V_\beta$ of $J_0$. Let $P\in\mathbb C[H_0]$
be an orthogonal projection supported on type $\gamma$, of rank
$r$ on $V_\gamma$, and let $S\in\mathbb C[J_0]$ be supported on
type $\beta$, of rank $s$ on $V_\beta$. On $\mathcal H$ each acts
on every multiplicity copy of its type. In particular $r,s$ are
ranks on the respective product-group irreducibles, not ranks on
$\mathcal H$ or on compact color carriers. We claim
\begin{equation}
 e^{(1-b)h(\mu,\xi)}{\rm tr}_{\mathcal H}(PS)
 \le r s (D_\gamma D_\beta)^{1-b}\exp O(t+n^{4/5}).
 \label{ten:compressed:eq:5}
\end{equation}
Zero ranks give zero overlap. For the other cases use the signed
tensor space on the non-tag sites, with $q$ even and $q$ odd colors.
Its global compact type $(\mu,\xi)$ is a tensor product of a
compact carrier and $\mathcal H$, by
\eqref{eq:signed-sector-decomposition}.
Restricted to row groups, multiplicities in the entire tensor
space are $\exp O(n^{4/5})$. Indeed, on each line the sum of
squared multiplicities is the commutant dimension. The signed
action on matrix units bounds that dimension by the number of
orbits, at most $(n+1)^{(2q)^2}$. There are $O(\sqrt n)$ lines
and $q=O(n^{1/16})$.
It follows that
$\operatorname{tr}_{\mathcal H}(PS)\le rs\exp O(n^{4/5})$.
We will interpolate this with a bound for the full row and column
type projections. Its proof keeps the compact carrier visible.

\label{ten:compressed:covariance-expansion}
We give the full covariance calculation, including the unitary multiplicity. This also explains why the comparison is valid for noncommuting densities.

If $n-t=0$, the untagged representation $\mathcal H$ is one dimensional, $h(\mu,\xi)=0$, and all local permutation types have dimension one. The overlap bound is immediate (a zero rank gives zero, and otherwise both ranks are one). We henceforth assume $n-t>0$.

Write $h=h(\mu,\xi)$ and let $\mathcal U$ be the irreducible representation of $U(q)\times U(q)$ of type $\tau=(\mu,\xi)$, with dimension $M$. On the global $\tau$-sector the tensor space is
\[
 \mathcal U\otimes\mathcal H.
\]
The symmetric-group action, and hence every row or column permutation projector, acts on $\mathcal H$ and as the identity on $\mathcal U$. For a tensor product $R$ of row densities and a tensor product $U$ of column densities, let $R_\tau,U_\tau$ be their diagonal blocks on this sector. Define the ordinary, unnormalized partial traces
\[
 R'=\operatorname{Tr}_{\mathcal U}R_\tau,
 \qquad U'=\operatorname{Tr}_{\mathcal U}U_\tau.
\]
Although $R$ and $U$ need not preserve the sector, their diagonal blocks are positive.

Let $\bar\sigma_R,\bar\sigma_C$ be the averages of the row and column densities over all corresponding line indices. First make the densities positive definite; the result for singular densities follows by approximation on their supports. Write
\[
 K_R=\rho_\tau(\bar\sigma_R^{-1/2}),\qquad
 K_C=\rho_\tau(\bar\sigma_C^{-1/2}).
\]
The highest-weight bound for a trace-one matrix gives
\[
 \rho_\tau(\bar\sigma_R)\le e^{-h}I_{\mathcal U},
 \qquad K_R^2\ge e^h I_{\mathcal U},
\]
and the same statements hold for $\bar\sigma_C$ and $K_C$. For every positive operator $A$ on $\mathcal U\otimes\mathcal H$,
\begin{equation}
\label{ten:compressed:partialtrace-positive}
 \operatorname{Tr}_{\mathcal U}
  \big((K_R\otimes I)A(K_R\otimes I)\big)
 \ge e^h\operatorname{Tr}_{\mathcal U}A.
\end{equation}
To verify the positive order, test a vector $v\in\mathcal H$. Its compression $A_v$ on $\mathcal U$ is positive, and the scalar on the left is
$\operatorname{Tr}(K_R A_vK_R)
 =\operatorname{Tr}(A_vK_R^2)
 \ge e^h\operatorname{Tr}A_v$.
In particular, the comparison does not require $K_R$ to commute with $A$.

Choose a parity-preserving matrix $Z_0$ with squared singular values
$(\mu_i/(n-t),\xi_j/(n-t))$ and put $Z=gZ_0h_0$, with $g,h_0$ independent Haar block unitaries. In the complete tensor space use
$X=\bar\sigma_R^{-1/2}Z\bar\sigma_C^{-1/2}$.
The product formula and Lemma~\ref{ten:missing-cells} give
\begin{equation}
\label{ten:compressed:covariance-upper}
 \operatorname{Tr}\big(RX^{\otimes(n-t)}
                U(X^*)^{\otimes(n-t)}\big)\le e^{O(t)}.
\end{equation}
Indeed the full-grid average of the one-site factors is
$\operatorname{Tr}(\bar\sigma_R X\bar\sigma_C X^*)=
 \operatorname{Tr}(ZZ^*)=1$.

Averaging independently over $g,h_0$ eliminates terms between inequivalent global unitary types by Schur orthogonality. Every surviving type contributes a nonnegative scalar. For the selected type $\tau$, the contribution is exactly
\begin{equation}
\label{ten:compressed:haar-identity}
 \frac{\|\rho_\tau(Z_0)\|_{\rm HS}^2}{M^2}
 \operatorname{Tr}_{\mathcal H}
 \left[
  \operatorname{Tr}_{\mathcal U}(K_R R_\tau K_R)
  \operatorname{Tr}_{\mathcal U}(K_C U_\tau K_C)
 \right].
\end{equation}
Here the carrier matrices are understood to be tensored with the identity. The formula follows twice from
$\int\rho(g)A\rho(g)^*\,dg=(\operatorname{Tr}A/M)I$;
expanding matrix entries on the multiplicity space gives the ordinary partial traces displayed above. The highest-weight monomial gives
$\|\rho_\tau(Z_0)\|_{\rm HS}^2\ge e^{-h}$.
Apply \eqref{ten:compressed:partialtrace-positive} to both positive blocks, and use positive trace comparison successively. Thus \eqref{ten:compressed:haar-identity} is at least
\[
 \frac{e^h}{M^2}\operatorname{Tr}_{\mathcal H}(R'U').
\]
Combining this with \eqref{ten:compressed:covariance-upper} yields
\begin{equation}
\label{ten:compressed:partialtrace-overlap}
 \operatorname{Tr}_{\mathcal H}(R'U')
       \le M^2e^{-h+O(t)}.
\end{equation}

We now specify the local density mixtures. In a row with $M_i>0$ untagged sites and permutation type $\gamma_i$, list all compatible polynomial type pairs $(\delta,\epsilon)$, each of length at most $q$, of total degree $M_i$, for which $\gamma_i$ occurs in $\delta*\epsilon^\top$. For each pair choose the parity-preserving trace-one matrix with combined spectrum $(\delta/M_i,\epsilon/M_i)$ and conjugate it by independent Haar unitaries in the two parity spaces. Its tensor power on the corresponding polynomial sector has, before averaging, highest-weight eigenvalue $e^{-h(\delta,\epsilon)}$. After averaging it is scalar there, with scalar at least that eigenvalue divided by the polynomial carrier dimension. It remains positive on every other sector.

By \eqref{ten:compressed:eq:4},
$h(\delta,\epsilon)\le\log D_{\gamma_i}+O(q^2\log n)$.
The carrier dimensions and the number of type pairs are both
$\exp(O(q^2\log n))$. Mix uniformly over this finite list of pairs. The resulting density average therefore dominates
$D_{\gamma_i}^{-1}e^{-O(q^2\log n)}$ times the full $\gamma_i$ projector on this row's signed tensor space: that projector is contained in the sum of the compatible polynomial sectors. A row with no untagged sites has scalar tensor power one and can use any parity-preserving density. Tensoring the independent row mixtures, and using the graded regrouping of sites, gives domination by
$D_\gamma^{-1}e^{-O(n^{4/5})}$ times the full row-type projector. Indeed there are $O(\sqrt n)$ lines and $q=O(n^{1/16})$, so all line errors fit in $O(n^{4/5})$. Construct the column mixture in exactly the same way with its local permutation types. It gives the corresponding factor with $D_\beta$. The row and column mixtures are independent. On the global sector these projectors have the form $I_{\mathcal U}\otimes P_\gamma$ and $I_{\mathcal U}\otimes P_\beta$. Their unnormalized partial traces are $M P_\gamma$ and $M P_\beta$. Independence of the two density mixtures and positivity therefore give the lower bound
\[
 \mathbb E\operatorname{Tr}(R'U')
 \ge \frac{M^2e^{-O(n^{4/5})}}{D_\gamma D_\beta}
                  \operatorname{Tr}_{\mathcal H}(P_\gamma P_\beta).
\]
The two factors $M$ cancel the $M^2$ in
\eqref{ten:compressed:partialtrace-overlap}, proving the full-type endpoint
\[
 \operatorname{Tr}_{\mathcal H}(P_\gamma P_\beta)
 \le D_\gamma D_\beta\,e^{-h+O(t+n^{4/5})}.
\]

Finally let $P,S$ have ranks $r,s>0$ in the respective irreducible symmetric-group spaces, as in the statement. Positive trace comparison bounds their overlap by the full-type endpoint. The multiplicity estimate supplies the independent bound
$\operatorname{Tr}(PS)\le rs\,e^{O(n^{4/5})}$.
For $0<b<1$, interpolate these two scalar upper bounds using
$\min(A,B)\le A^bB^{1-b}$; since $(rs)^b\le rs$, the result is
\[
 e^{(1-b)h}\operatorname{Tr}_{\mathcal H}(PS)
 \le rs(D_\gamma D_\beta)^{1-b}e^{O(t+n^{4/5})},
\]
which is the stated overlap estimate. The case of a zero rank is immediate.


\subsection{Counting the tag layouts}
Spectrally expanding \(A_0,B_0\) into positive weighted projections and summing \eqref{ten:compressed:eq:5} using \eqref{ten:compressed:eq:3} (numbers of types absorbed), we conclude
\[
 \begin{split}
 D_\lambda {\rm tr}_\lambda AB
 &\le \exp\{b h(\mu,\xi)+c t\log n+O(t+n^{1-c/6})\}\\
 &\qquad{}\cdot \sum_{|T|=t}\exp\{-\sum_i r_i\log(m/r_i)-\sum_j s_j\log(m'/s_j)\}.
 \end{split}
\]
Here per table we used \(\binom nt t!(|H_0||J_0|)/(|H||J|)\le e^{O(t)}\), by falling factorial bounds. Group the tables by margins (subexponentially many, \(\exp O(n^{4/5})\)). With given margins their number is at most
\[
 \exp\{\sum_i r_i\log(m/r_i)+\sum_j s_j\log(m'/s_j)-t\log(n/t)+O(t)\}.
\]
Indeed order the \(t\) entries arbitrarily, use multinomials for the two coordinate sequences separately and divide by \(t!\). This gives precisely the desired budget with absorbable losses.

\subsection{Closing the exponent recursion}
Fix \(b\), then \(c\), and then \(\eta\) in the ranges required above, in particular with \(\eta<c/6\). The raw tag calculation gives
\(L_n\ge-n^{1-c/2}/e\). Fix the theorem's allowance \(C\ge1/e\) now, so
\(L_n+C n^{1-c/2}\ge0\) at every size and for every diagram.

For a dense block that is not annihilated, write \(W=\log D_\lambda\). The bounds on \(t\) and \(W\), together with \(\log n=d\log2\) and \(\eta<c/6\), give
\(t+n^{1-c/6}\le K_0W/d\) above a fixed size. Consequently the preceding trace and layout estimates, with this \(C\) fixed, give
\[
 \begin{aligned}
 D_\lambda {\rm tr}|\rho_\lambda(B_n)|^p
 &\le \exp\{L_n(\lambda)+K W/d\},\\
 p&=\max(p_{\lfloor d/2\rfloor},p_{\lceil d/2\rceil}),
 \end{aligned}
\]
whenever the child bounds hold with exponents at least two. Here \(K\) is independent of \(p\): the trace inequality has unit constant, and the coefficients in the remaining errors above depend only on the fixed budget constants and allowance. Increase the size cutoff so that \(K/d\le1/4\). Since \(L_n\le W/2\), the last display makes every singular value \(s\) satisfy \(s^p\le e^{-W/4}\). Hence
\[
 \begin{aligned}
 D_\lambda {\rm tr}|\rho_\lambda(B_n)|^{p(1+A/d)}
 &\le \exp\{L_n(\lambda)+(K-A/4)W/d\}\\
 &\le e^{L_n(\lambda)}
 \end{aligned}
\]
when \(A\ge4K\). Choose an integer cutoff \(d_0\ge1\) large enough for this argument, the minimizing-part and sparse estimates, and all child-size thresholds used above. Both \(A\) and \(d_0\) have now been fixed independently of the base power.

Choose one real \(P_*\ge2\) above the fixed sparse requirement, including the stronger sparse decay, and large enough for every nontrivial block with \(1\le d\le d_0\). This is possible by Lemma~\ref{lem:finitegap}: there are finitely many such blocks, their sweep norms are strictly below one, and
\[
 \begin{aligned}
 D_\lambda {\rm tr}|\rho_\lambda(B_n)|^{P_*}
 &\le D_\lambda^2\|\rho_\lambda(B_n)\|^{P_*}\\
 &\le1
 \end{aligned}
\]
once \(P_*\) is large enough. Trivial blocks have weighted trace one and are covered by the already fixed nonnegative allowance. Set
\[
 \begin{aligned}
 p_d&=P_*\quad(1\le d\le d_0),\\
 p_d&=(1+A/d)\max(p_{\lfloor d/2\rfloor},p_{\lceil d/2\rceil})
       \quad(d>d_0).
 \end{aligned}
\]
Every \(p_d\) is at least \(P_*\), so the sparse estimate remains available. Along a rounded-halving branch above the cutoff, a parent \(u\) and child \(v\) satisfy \(u-1\ge2(v-1)\). Thus the sum of \(1/u\) along that branch is at most \(2/d_0\), and
\(p_d\le P_*\exp(2A/d_0)\). This proves \eqref{ten:compressed:eq:1} with bounded exponents.

After a sufficiently large absolute constant number of sweeps, Plancherel now gives vanishing squared relative \(L^2\) distance to uniform: for dense types \eqref{ten:compressed:eq:1} gives individual norm to bounded power \(\le e^{-\Omega(\log D_\lambda)}\), dominating both \(D_\lambda^2\) and the partition count; for sparse types use \eqref{ten:compressed:eq:2} and \(D_\lambda\le n^k\), with at most \(2^k\) diagrams per level. This implies the total variation bound, independent of the starting order. The support obstruction is Lemma~\ref{lem:support}.
